\documentclass[11pt,a4paper]{article}

\usepackage{amsmath}
\usepackage{amssymb}
\usepackage{graphicx}
\usepackage{geometry}
\geometry{margin=2.5cm}

\title{Local Symmetry--Adapted Internal Coordinates:\
A Topological and Group--Free Construction}
\author{}
\date{}

\begin{document}
\maketitle

\begin{abstract}
We present a fully automatic strategy for the construction of non--redundant,
local symmetry--adapted internal coordinates for molecular systems.
The approach is primarily topological and does not require the explicit
assignment of a global point group, making it robust to weak distortions,
conformational variability, and large flexible systems.
Equivalence relations between internal coordinates are inferred from atomic
classes, local environments, and ring connectivity, allowing the generation of
linearly independent collective coordinates suitable for vibrational analysis,
reduced--dimensional modeling, and large molecular systems.

A hybrid option is introduced: a topology/geometry--driven symmetry basis is
built first and redundancies are removed optionally by a final pruning step
based on the metric (G) or SVD. This preserves continuity of the coordinate
basis while ensuring independence when local degeneracies occur, and provides
a reproducible, model--agnostic route to symmetry--aware coordinates without
global group assignments.
\end{abstract}

\section{Motivation and Scope}

Internal coordinates are the natural language for describing molecular
vibrations, structural deformations, and reaction pathways.
However, their direct use often leads to large sets of redundant variables,
especially in molecules with local symmetries, equivalent substituents,
or condensed ring systems.

Traditional symmetry--adapted approaches rely on the explicit identification
of the molecular point group and the subsequent projection onto irreducible
representations.
While rigorous, this strategy becomes cumbersome or ambiguous for large,
flexible, or weakly distorted systems.

Natural internal coordinates in the sense of Pulay provide a powerful
chemically intuitive framework and are implemented in codes such as Molrp.\cite{pulay_natural,molrp}
However, the typical construction starts from redundant primitive sets and
requires careful handling of linear dependencies. At the opposite extreme,
fully delocalized non--redundant coordinates (e.g.\ Baker's delocalized
internal coordinates) can lose direct physical interpretability.\cite{baker_deloc}
The present work aims at a middle ground: locally symmetry--adapted,
topology--driven coordinates that remain interpretable while avoiding the
pathologies of redundant primitives and the opacity of highly delocalized bases.

The goal of the present work is to construct \emph{local symmetry--adapted
internal coordinates} without invoking a global point group, relying instead on
topological equivalence, local environments, graph--based analysis, and
geometry--assisted local symmetry recognition. This yields a robust and
fully automated pipeline that is stable under small perturbations and suitable
for high--throughput workflows.

\subsection{Key Contributions}
\begin{itemize}
  \item Introduction of \emph{continuous, synthon--based effective atomic numbers}
        ($Z_{\\mathrm{eff}}$) as the primary criterion for local equivalence
        classes, enabling robust symmetry inference without a global point group.
  \item A topology--first, geometry--aware construction of local symmetry bases
        that remains interpretable while avoiding redundancy and discontinuities.
  \item Systematic cyclic coordinate definitions for rings (breathing, CVB,
        cyclic torsions) with deterministic handling of fused/condensed
        ring dependencies, plus inter--ring butterfly/hinge coordinates.
  \item Integration of TRIC--style inter--fragment coordinates into the same
        symmetry pipeline for disconnected systems.
  \item Optional per--block pruning (SVD or G--metric) that removes residual
        dependencies without reordering the underlying local basis.
\end{itemize}

\section{General Philosophy}

The methodology is based on the following principles:

\begin{itemize}
  \item Only internal coordinates of the same type are combined
        (bond lengths, valence angles, out--of--plane angles, torsions,
        linear--bend components).
  \item X--H bond lengths are never mixed with non--hydrogen bonds.
  \item Geminal X--H coordinates may be combined, while vicinal or remote
        combinations are excluded.
  \item Atomic equivalence is inferred from topological and chemical similarity,
        not necessarily exact symmetry.
  \item Linear combinations are constructed to remove redundancies while
        preserving physical interpretability.
  \item Non--cyclic torsions are reduced to one dihedral per bond using
        substituent priority derived from effective atomic numbers.
\end{itemize}

The resulting coordinates are suitable for vibrational Hamiltonians,
VSCF/VPT2 treatments, and reduced--dimensional models.

\section{Atomic Classes, Zeff, and Local Geometry}

Atoms are first grouped into equivalence classes based on:
\begin{itemize}
  \item Effective atomic number $Z_{\mathrm{eff}}$ with a tunable threshold,
  \item Coordination number and discrete topology descriptors,
  \item Local bonding pattern, aromatic/ring membership, and chemical environment.
\end{itemize}

This yields discrete equivalence classes that are stable under small
geometrical deformations.

In addition, local geometry around each center is recognized when possible
(e.g.\ tetrahedral vs square--planar for coordination 4; TBP vs square--pyramidal
for coordination 5; octahedral vs trigonal--prismatic for coordination 6; and
square--antiprismatic, dodecahedral, or cubic for coordination 8).
This allows local symmetry bases to be chosen consistently with the actual
coordination geometry.

\section{Primitive Internal Coordinates}

Primitive internal coordinates are generated:
\begin{itemize}
  \item Bond stretches,
  \item Valence angles,
  \item Wilson--type out--of--plane angles,
  \item Proper torsional angles (no impropers),
  \item Linear--bend components for near--linear centers.
\end{itemize}

Each primitive coordinate is associated with a signature encoding the
equivalence classes of the atoms involved, the coordinate type,
and (optionally) a coarse angular bin.

\section{Grouping of Equivalent Primitives}

Primitive coordinates sharing identical signatures are grouped together.
These groups represent sets of coordinates that are physically equivalent or
nearly equivalent (e.g.\ the three Y--X--Y angles in an XY$_3$ moiety).

For each group, a deterministic local basis is constructed:
\begin{itemize}
  \item A symmetric sum coordinate,
  \item Orthogonal difference coordinates (Gram--Schmidt on simple differences).
\end{itemize}

When local geometry and class patterns allow, symmetry--specific bases are used
(e.g.\ XY$_3$ with $C_{3v}$; XY$_2$Z with $C_{2v}$; XY$_4$ with $T_d$ or
square--planar patterns). Otherwise, the generic basis is used.

\section{Ring Detection and Canonicalization}

Rings are identified purely from molecular connectivity using graph traversal.
Each ring is represented as an ordered list of atoms.

To avoid duplicates and orientation ambiguities, each ring is canonicalized
by selecting a unique cyclic permutation and direction.
This allows consistent identification of equivalent rings and condensed
ring systems.

\section{Ring Clusters and Redundancy Reduction}

Rings sharing one or more atoms are grouped into ring clusters.
For a cluster containing $M$ rings, only $M-1$ independent collective ring
coordinates are retained.
This reflects the intrinsic linear dependency present in fused or condensed
ring systems.

The discarded coordinate is chosen deterministically to ensure reproducibility.

\section{Collective Ring Coordinates}

For each retained ring, the following collective coordinates are constructed:

\begin{itemize}
  \item Ring breathing modes (uniform expansion/contraction),
  \item Cyclic valence bending (CVB) coordinates,
  \item Cyclic torsional coordinates.
\end{itemize}

These coordinates are defined as linear combinations of primitive internal
coordinates associated with the atoms of the ring and are ordered according to
the canonical ring traversal.
Butterfly coordinates for fused rings sharing a bond are treated as dedicated
blocks to capture out--of--plane collective motion. Condensed ring clusters
also generate hinge--like inter--ring torsional coordinates, with one
redundant combination removed per cluster (the $M-1$ rule).

\section{Inter--Fragment Coordinates (TRIC/geomeTRIC)}

For disconnected systems, inter--fragment coordinates are added to preserve
rigid relative placement between fragments and to avoid singular transformations.
We follow the TRIC philosophy as implemented in geomeTRIC: translations are
defined using fragment centroids, and rotations are defined through an
exponential map of the quaternion describing the relative orientation between
fragment frames. These six coordinates (three translational and three rotational
per non--reference fragment) are equivalent in spirit to the inter--fragment
coordinates used in geomeTRIC, and integrate seamlessly with the local
symmetry pipeline described above.\cite{wang2016tric}

\section{Non--Redundancy and Hybrid Pruning}

The final set of collective coordinates is non--redundant by construction.
Redundancy is eliminated locally:
\begin{itemize}
  \item within groups of equivalent primitives,
  \item within sets of equivalent rings,
  \item within condensed ring clusters.
\end{itemize}

A hybrid option is provided:
\begin{enumerate}
  \item Construct a topology/geometry--driven symmetry basis $U_{\mathrm{topo}}$,
  \item Optionally prune residual dependencies using either
        (a) SVD of $U_{\mathrm{topo}}$ or
        (b) a metric--based pruning on $G_q = U_{\mathrm{topo}}^T G U_{\mathrm{topo}}$.
\end{enumerate}

This preserves continuity and interpretability of the coordinate basis while
ensuring independence when local degeneracies occur.

In this sense, the proposed construction complements Pulay--type natural
coordinates (as in Molrp)\cite{pulay_natural,molrp} by removing redundancies through topology and local
symmetry, while retaining more local character than fully delocalized
non--redundant schemes (e.g.\ Baker).\cite{baker_deloc} The result is a robust, scalable set of
coordinates that balances interpretability and numerical stability.

\section{Global Symmetrization (Optional)}

Once a non--redundant $U$ is constructed, an optional global symmetrization can
be applied using the point--group symmetry of the entire molecule. The workflow
assumes that the geometry is centered and oriented along its principal axes.
Symmetry operations are detected by matching the oriented coordinates under a
set of candidate operations (including $E$, inversion, mirror planes, proper
rotations $C_n$ up to a user--defined $n_{\max}$, and improper rotations $S_n$).
The corresponding atomic permutations are propagated to the primitive
coordinates, and the columns of $U$ are projected onto the symmetric subspace.

When requested, the procedure can retain only the totally symmetric (A$_1$)
coordinates, which is particularly useful for constrained or symmetry--preserving
geometry optimizations. Symmetrization is applied \emph{per coordinate type}
to avoid mixing stretches, bends, torsions, and other primitive families.\cite{cotton_group}

Additional global options include: (i) separate tolerances for hydrogen atoms,
(ii) heavy--atom orientation for robust principal axes, (iii) optional strict
filters on symmetry operation deviations, (iv) isotopic sensitivity
(treating isotopes as equivalent or distinct), (v) a relative tolerance term
that scales with distance from the origin, and (vi) an optional inertia--based
reduction of the maximum $n$ to avoid unnecessary testing of high--order axes
for strongly asymmetric tops. The global symmetry report can include the
detected operations, a point--group label, generators, and deviation diagnostics
(mean/max), and may optionally include a confidence score and lightweight
profiling statistics (number of operations tested/accepted).
\medskip
\noindent\textbf{Additional controls.} Matching can be scaled by an explicit
radius (for uniform tolerance across system sizes), the radial pre--filter may
be disabled for debugging, and candidate operations can be restricted to
subgroup families (e.g.\ $C_n$, $D_n$, or polyhedral) when desired.
\medskip
\noindent\textbf{Implementation detail.} To reduce cost, candidate operations are filtered with
lightweight invariants (principal moments), and atomic matching is accelerated by
symbol grouping, radial pre--filters, and early exits when no compatible matches
remain. These steps preserve exactness within the chosen tolerances while improving
scalability.
\medskip
\noindent\textbf{Reporting.} The text report can optionally include a compact global
symmetry summary with confidence estimates and lightweight profiling counters
(operations tested/accepted), aiding diagnostics in large workflows.

\noindent\textbf{Example.} For symmetry--preserving geometry optimizations, one may
construct $U$ and then retain only A$_1$ coordinates, optionally recording symmetry
labels and diagnostics in the report.
For isotopically substituted systems, one may treat isotopes as distinct or
equivalent, and optionally enforce a strict deviation cutoff for accepted
symmetry operations.

\section{Implementation}

The entire framework is implemented in a strict procedural style,
compatible with legacy Fortran workflows and modern scripting interfaces.
No global symmetry assignment is required.

The implementation is organized in sequential steps and is designed to be
robust, deterministic, and extensible:
\begin{enumerate}
  \item Atomic classification (including $Z_{\mathrm{eff}}$ thresholding),
  \item Primitive coordinate generation,
  \item Signature construction and grouping,
  \item Local geometry recognition (when applicable),
  \item Ring detection and canonicalization,
  \item Ring cluster analysis,
  \item Construction of collective coordinates,
  \item Optional SVD/G pruning of residual dependencies.
\end{enumerate}

Pattern reports can be produced to summarize local symmetries and
group assignments, enabling transparent diagnostics and reproducibility.

\subsection{Pipeline Overview (Pseudocode)}

\begin{verbatim}
1. Build topology (bonds, rings) and atomic descriptors
2. Compute Zeff-based atomic equivalence classes
3. Generate primitive internal coordinates
4. Group primitives by signatures (type, classes, ring tag, angle bin)
5. Recognize local geometry (coord. 4/5/6/8) where applicable
6. Build local symmetry bases (or generic sum/diff fallback)
7. Add cyclic ring coordinates (breathing, CVB, torsions, butterfly)
8. Reduce fused-ring redundancies (M-1 rule)
9. Optionally prune per block (SVD or G)
10. Output U and pattern report
\end{verbatim}

\subsection{Limitations and Scope}

The local geometry recognition step is designed to be robust, but ambiguous
cases can arise (e.g.\ between closely related coordination geometries). In such
cases the pipeline falls back to a generic, deterministic basis to preserve
continuity. The present framework is intended for molecular systems where local
coordination environments can be classified reliably; highly fluxional systems
may benefit from tighter thresholds or trajectory--based smoothing.

\subsection{Typical Parameters}

\begin{center}
\begin{tabular}{ll}
\hline
Parameter & Typical value \\\\
\hline
$Z_{\\mathrm{eff}}$ threshold & 0.05 \\\\
Geometry match tolerance & 12$^{\\circ}$ \\\\
Pruning mode & SVD (per block) \\\\
\hline
\end{tabular}
\end{center}

\section{Outlook}

The present framework provides a robust foundation for:
\begin{itemize}
  \item automatic generation of reduced internal coordinate sets,
  \item vibrational calculations on large and flexible systems,
  \item data--driven or fragment--based vibrational models.
\end{itemize}

Future developments will include refined local symmetry libraries,
improved geometry matching, coupling with force--constant generation,
and integration into composite quantum--chemical workflows. The present
implementation already provides an innovative, geometry--aware and
topology--driven alternative to traditional symmetry projection methods.

An additional advantage is the continuity of the coordinate basis along
reaction paths or molecular dynamics trajectories: local equivalence classes
and cyclic coordinate ordering provide stable coordinate identities, while the
optional pruning step only removes residual dependencies without reordering
the underlying basis.

\section{Computational Cost}

The pipeline is dominated by primitive generation and local grouping.
Per--block pruning scales modestly with group size, and cyclic ring operations
scale linearly with the number of ring primitives. In practice, the approach
is suitable for large systems and can be accelerated through caching of
primitive values and B--matrices.

\appendix
\section{Implementation Overview and Routine Dependencies}

This appendix provides a structured overview of the main routines composing
the local symmetry--adapted internal coordinate framework, together with their
primary dependencies.
The goal is to clarify data flow, logical separation of tasks, and code
maintainability.

\subsection{Input and Preprocessing}

\begin{itemize}
  \item \texttt{READ\_XYZIN} \\
  Reads molecular geometry and optional metadata from an enriched XYZ input
  file (Merlino format).
  \begin{itemize}
    \item Outputs: atomic numbers, Cartesian coordinates
  \end{itemize}

  \item \texttt{READ\_CONNECTIVITY} \\
  Reads or constructs molecular connectivity (bond list).
  \begin{itemize}
    \item Outputs: bond pairs
  \end{itemize}
\end{itemize}

\subsection{Atomic Classification}

\begin{itemize}
  \item \texttt{BUILD\_ATOM\_CLASSES} \\
  Assigns atoms to equivalence classes based on $Z_{\mathrm{eff}}$ and local
  connectivity descriptors.
  \begin{itemize}
    \item Depends on: connectivity
    \item Outputs: atomic class labels
  \end{itemize}
\end{itemize}

\subsection{Primitive Internal Coordinates}

\begin{itemize}
  \item \texttt{BUILD\_PRIMITIVE\_COORDS} \\
  Generates the full set of primitive internal coordinates:
  bond stretches, valence angles, Wilson out--of--plane angles, proper
  torsions (no impropers), and linear--bend components for near--linear
  centers.
  \begin{itemize}
    \item Depends on: geometry, connectivity
    \item Outputs: primitive coordinate list and atom indices
  \end{itemize}
\end{itemize}

\subsection{Signature Construction and Grouping}

\begin{itemize}
  \item \texttt{BUILD\_SIGNATURES} \\
  Constructs topological signatures for primitive coordinates using atomic
  classes, coordinate types, and coarse angular bins.

  \item \texttt{GROUP\_PRIMITIVES} \\
  Groups primitives sharing identical signatures into equivalence sets.
  \begin{itemize}
    \item Depends on: \texttt{BUILD\_SIGNATURES}
    \item Outputs: primitive groups
  \end{itemize}
\end{itemize}

\subsection{Local Geometry Recognition}

\begin{itemize}
  \item \texttt{MATCH\_LOCAL\_GEOMETRY} \\
  Matches local coordination environments to geometry templates for
  coordination 4/5/6/8 where applicable.
\end{itemize}

\subsection{Ring Analysis}

\begin{itemize}
  \item \texttt{FIND\_RINGS} \\
  Identifies rings from molecular connectivity using graph traversal.
  \begin{itemize}
    \item Depends on: connectivity
    \item Outputs: raw ring lists
  \end{itemize}

  \item \texttt{CANONICALIZE\_RING} \\
  Produces a canonical representation of a ring by cyclic permutation and
  direction selection.

  \item \texttt{NORMALIZE\_RINGS} \\
  Applies canonicalization to all detected rings.

  \item \texttt{REMOVE\_DUPLICATE\_RINGS} \\
  Eliminates identical rings after canonicalization.
\end{itemize}

\subsection{Ring Clustering}

\begin{itemize}
  \item \texttt{BUILD\_RING\_ATOM\_MAP} \\
  Builds atom--to--ring incidence lists.

  \item \texttt{SHARE\_ATOM} \\
  Logical function testing whether two rings share at least one atom.

  \item \texttt{BUILD\_RING\_CLUSTERS} \\
  Groups rings into clusters based on shared atoms.

  \item \texttt{REDUCE\_RING\_CLUSTER\_N1} \\
  Applies the $M-1$ reduction rule to each ring cluster to remove linear
  dependencies.
\end{itemize}

\subsection{Collective Coordinate Construction}

\begin{itemize}
  \item \texttt{BUILD\_COLLECTIVE\_FROM\_GROUP} \\
  Constructs collective coordinates from groups of equivalent primitives.

  \item \texttt{BUILD\_RING\_BREATHING} \\
  Builds ring breathing coordinates (cyclic bond sums).

  \item \texttt{BUILD\_CVB\_FOR\_RING} \\
  Builds cyclic valence bending coordinates for a ring.

  \item \texttt{BUILD\_CT\_FOR\_RING} \\
  Builds cyclic torsional coordinates for a ring.

  \item \texttt{BUILD\_BUTTERFLY\_FOR\_FUSED} \\
  Builds butterfly coordinates for fused rings.

  \item \texttt{BUILD\_HINGE\_FOR\_CLUSTERS} \\
  Builds hinge--like inter--ring torsional coordinates for condensed ring
  clusters (with $M-1$ reduction).

  \item \texttt{BUILD\_ALL\_COLLECTIVE\_COORDS} \\
  High--level driver assembling all collective coordinates:
  local groups, ring breathing modes, CVB, cyclic torsions, butterfly/hinge
  coordinates, and cluster reductions.
\end{itemize}

\subsection{Optional Pruning}

\begin{itemize}
  \item \texttt{PRUNE\_U\_SVD} \\
  Removes residual dependencies by SVD of the assembled $U$.

  \item \texttt{PRUNE\_U\_G} \\
  Removes residual dependencies using $G_q = U^T G U$.
\end{itemize}

\subsection{Main Driver}

\begin{itemize}
  \item \texttt{MAIN\_LOCAL\_COORDS} \\
  Sequential driver routine coordinating all steps of the workflow.
\end{itemize}

\subsection{Design Philosophy}

The code is strictly procedural and compatible with Fortran~77 / old--style
Fortran~90.
All dependencies are explicit through subroutine arguments; no global state,
modules, or implicit interfaces are used.
This design favors transparency, portability, and long--term maintainability.

\section{Local Pattern Reference}

Representative local patterns inferred from equivalence classes and geometry:
\begin{itemize}
  \item Coord. 3: XY$_3$ (C$_{3v}$), XY$_2$Z (C$_{2v}$), XYZ (C$_1$)
  \item Coord. 4: XY$_4$ (T$_d$), square--planar (D$_{4h}$), XY$_3$Z (C$_{3v}$),
        XY$_2$Z$_2$ (C$_{2v}$/D$_{2h}$)
  \item Coord. 5: TBP vs square--pyramidal (geometry matched)
  \item Coord. 6: octahedral vs trigonal--prismatic
  \item Coord. 8: square--antiprismatic, dodecahedral, cubic
\end{itemize}

\begin{thebibliography}{9}
\bibitem{wang2016tric}
Wang, L.-P.; Song, C.
\emph{Geometry Optimization Made Simple with Translation and Rotation Coordinates.}
J. Chem. Phys. \textbf{2016}, 144, 214108.
\bibitem{pulay_natural}
Pulay, P.
\emph{Ab initio calculation of force constants and equilibrium geometries in polyatomic molecules.}
Mol. Phys. \textbf{1969}, 17, 197--204.
\bibitem{molrp}
Rappoport, D.; Furche, F.
\emph{Molecular vibration analysis with natural internal coordinates.}
Mol. Phys. \textbf{2010}, 108, 531--552.
\bibitem{baker_deloc}
Baker, J.
\emph{Techniques for geometry optimization: A comparison of Cartesian and internal coordinates.}
J. Comput. Chem. \textbf{1993}, 14, 1085--1100.
\bibitem{cotton_group}
Cotton, F. A.
\emph{Chemical Applications of Group Theory}, 3rd ed.;
Wiley: New York, 1990.
\end{thebibliography}

\end{document}
